\section{模型评价与推广}

\subsection{优点}

交会定位把示向度误差写成闭楔形而不是点估计，直径与 Jung 半径分开计算，避免用 $\mathrm{diam}/2$ 充当光学半径。
等边三角形演示给出直径约 $1.000$、Jung 半径约 $0.577$，直径圆不能覆盖；两点楔形演示给出直径约 $9.877$、包围半径约 $4.939$，直径圆恰好覆盖。
两种情形并列，说明覆盖性按顶点几何判定，而不是写死布尔值。
问题2的 $J(q)$ 只承担排序，演示 $J^\ast\approx 43.03\,\mathrm{m}$ 明确大于 $20\,\mathrm{m}$，不签发清除证书，与后文光学网格的分工清楚。
问题3、问题4把字典序目标落到可检查的证书：发现点集上的频道判空，加上连续候选区上的 $20\,\mathrm{m}$ 光学覆盖。
HEX37 的覆盖半径 $\delta\approx 461.880\,\mathrm{m}$、半平面间隙约 $38.120\,\mathrm{m}$ 与第 $4$ 圈排除间隙约 $9.401\,\mathrm{m}$ 有显式不等式，不是演练碰巧全清。
累计楔形裁剪保留顶点顺序，离线单场时间从约 $11\,\mathrm{s}$ 降到约 $0.1$--$0.4\,\mathrm{s}$，使三十种子配对成为可重复实验而不是一次不可复现的运行。

相对仅用到达方向做统计定位的工作\upcite{weiss_2005_dpd,penna_2013_doa}，本文必须在未知 $N$、未知 $r_i$、单机单频道与硬性时间账本下同时完成清除。相对覆盖路径规划综述中的格子扫掠\upcite{galceran_2013_cpp,cabreira_2019_uavcpp}，格点还要进入未知 $180^\circ$ 半平面。相对定向越野在时间预算下挑选收益顶点\upcite{vansteenwegen_2011_op,gunawan_2016_op}，本题漏清不可用超时换取，完备性证书必须放在目标函数之外作为约束。这些约束使“证书优先、启发式插入其次”的分层比单层黑箱搜索更符合题面。问题3主线 GREEDY\_FAST 在非正式十局中全清 $10/10$、平均 $T/K=341\,\mathrm{s}$，相对 GREEDY 的 $390\,\mathrm{s}$ 下降来自更少走廊而不是放弃全清；问题4离线配对上 HEX37 相对 SQUARE81 的 $T/K$ 下降 $36.267\%$，全程 $30/30$ 证书。回归测试 $168$ 项通过。

\subsection{缺点与未完成项}

正式测试未执行，表1空缺。
主线策略没有在官方正式模块上验证，截止之后不能补做。
演练局数少且非配对，不能把 $341\,\mathrm{s}$ 或 $822.38\,\mathrm{s}$ 说成期望值，也不能把离线 $836.91\,\mathrm{s}$ 或 PREFIX\_A 的 $683.09\,\mathrm{s}$ 写入正式成绩。
把这四组数字平均成一个现场成绩，会把演练、离线几何模型与正式测试三条口径混在一起，比算错直径更严重。

问题3的七点 $P_3$ 对 $r_i\in[1000,1200)$ 不是单点可测的严格证书，完备性更多依赖巡回之后的局部走廊。
问题2的 $J(q)$ 用边界采样，连续集合上的最坏半径可能被低估；开发集上真值后验半径最大仍有 $44.24\,\mathrm{m}$，超过光学半径。
HEX37 未证明真子集仍覆盖一切 $(g,d)$，减点必须另证。
光学方格在定位域很大时点数上升，全定向、后验长期张得很开的场景仍会走较多兜底。
SQUARE81 演练一场清除尝试 $223$ 次才成功 $12$ 次，说明兜底代价可以很高。
$N=16$ 的计数证书尚未改成提前停机，四条空频道仍会付出覆盖测向。

实现仍可能在覆盖点上对已 $\mathtt{detected}$ 的频道重复测量（证书模式），采样种子上该路径被“每站立刻处理待清除”掩盖，但代码路径存在。墙钟侧的 UIA 启动偶发误超时，与虚拟时间不是同一问题，却可能吃掉正式测试额度。离线几何模型不经过 HTTP 接口，不能替代模拟器正式测试。数值包络 $1.005^\circ$ 是量化约定，不是题面参数。

\subsection{推广}

同一套楔形交与光学方格可以迁到多机器人：证书层仍要保证半平面与接收半径，任务分配应使用联合目标，避免一辆车跑完覆盖、另一辆闲置。
三维到达角定位的误差界有闭合形式讨论\upcite{zhao_2013_placement,li_2020_crlb}，若测向机提供俯仰角，闭楔形应改成闭圆锥交。
未知个数的射频源搜索与定向越野的收益顶点类似\upcite{gunawan_2016_op}，但本题漏检不可用超时换取，推广时仍应把完备性证书放在目标函数之外作为约束。
传感器网络中的检测、分类与跟踪层次\upcite{li_2002_tracking}在多机、多频道下可以重新展开，本题被单机单频道压成一条回路，并不否定分层本身。

若接收半径下界提高，HEX37 的步长可以放大；若定向扇窄于 $180^\circ$，现有半平面证明失效，必须改点集。任何把示向度误差从一度放宽的场景，Jung 半径与光学方格边长都要重算，不能沿用 $28\,\mathrm{m}$。若正式测试额度允许且截止未到，应先用演练入口核对证书，再独立启动正式模块，不得把演练日志改名。本文把完备性证书放在启发式加速之前，这一原则在更换点集、更换插入阈值或更换光学边长时仍然适用。

写作时正式测试未执行，因此本文能够确定的是算法对象与证书，而不是竞赛名次。
能够确定的对象包括：闭楔形交集、直径与 Jung 半径的分离、第二点排序指标 $J$ 的非证书地位、$P_3$ 七点发现网格、SQUARE81 与 HEX37 的覆盖不等式、边长 $28\,\mathrm{m}$ 光学方格的半对角线条件，以及 GREEDY\_FAST 与 PREFIX\_A 在非正式口径下的筛选结果。
不能确定的对象包括：正式三次的 $K$、$T/K$ 与墙钟，HEX37 真子集的覆盖性，以及把演练均值外推到官方场景的期望值。
把能确定的写成命题，把不能确定的留空，比用演练数字填满表1更符合题面。
非正式口径下可以引用的数字已经冻结：问题3主线十局 $341\,\mathrm{s}$，问题4离线配对 HEX37 为 $836.91\,\mathrm{s}$，PREFIX\_A 离线均值为 $683.09\,\mathrm{s}$，演练两场为 $822.38\,\mathrm{s}$ 与 $930.28\,\mathrm{s}$。
SQUARE81 修复对照把全清率从 $13/30$ 拉到 $30/30$，平均 $T/K$ 从 $2298.44\,\mathrm{s}$ 降到 $1313.16\,\mathrm{s}$；留出种子 $1000$--$1029$ 仍为 $30/30$，均值 $1376.27\,\mathrm{s}$。
这些数字分属演练与离线几何模型，禁止平均成一个现场成绩，也禁止把修复前含漏清的 $2298.44\,\mathrm{s}$ 单独拿来和修复后比快慢。
问题2开发集上 $J^\ast$ 全体高于 $20\,\mathrm{m}$、一级失败计数全为零，只说明排序模块没有签发虚假证书，不能外推到正式测试。

从方法上看，本文没有把四问压成一个黑箱优化器。
问题1是凸集几何，问题2是外包络上的排序，问题3、问题4是带证书的在线闭环。
分层的代价是实现变长：覆盖点集、累计裁剪、光学方格、插入寻路与界面自动化必须分开测试。
分层的收益是失败可定位：旧版漏清可以归到预算停机和观测不累计，而不是归到网格点数不够。
HEX37 把发现行程缩短之后，光学兜底次数几乎不变（约 $1.07$ 对 $1.17$），也说明清除层已经从发现层解耦。
若将来要再加速，应先证更稀的点集或更早的计数判空，而不是把 SAFE 走廊改成随机跳点。
$N=16$ 时十六个不同频道一旦都出现正观测，其余四频道可由计数判空，这是计数证书而不是几何证书；当前实现仍走完覆盖点，以免把两种证书混成提前停机。
附录只列出与命题直接对应的片段，完整源文件放在支撑材料，以免把实现细节误写成新的题设参数。
参数表 $28$ 项、正式额度每问三次、截止时间 $2026$ 年 $9$ 月 $13$ 日 $17$:$30$，均按题面冻结，正文不另改数。
